\begin{frame}[standout]

    Focus sur Apache Cassandra

    \note{
        Cassandra est l'exemple central de la famille orientée colonne
        dans ce cours. L'objectif est de relier son architecture aux besoins
        de scalabilité, de disponibilité et de débit.
    }

\end{frame}

\begin{frame}{Qu'est-ce que Cassandra ?}

    %    Au départ, Cassandra est un projet développé par Facebook avec pour fonctionnalité la recherche de messages dans la boîte de réception de Facebook.
    %    En juillet 2008, Facebook publie le code source en tant que projet open source sur Google Code.
    %    Puis, en mars 2009, il est incubé au sein de la fondation Apache Software Foundation, et devient un projet de premier rang en février 2010.

    \begin{figure}[htb]
        \resizebox{\textwidth}{!}{
        \begin{tikzpicture}
            \tikzstyle{descript}=[text=black,align=center,minimum height=1.8cm,align=center,outer sep=0pt]
            \tikzstyle{activity}=[align=center,outer sep=1pt]

            %% Coordinates
            \coordinate (O) at (-1,0); % Origin
            \coordinate (F) at (14,0); % End
            \coordinate (P1) at (0,0);
            \coordinate (P2) at (2,0);
            \coordinate (P3) at (12,0);
            \coordinate (P4) at (12,0);

            %% Events
            \coordinate (E1) at (1.8,0); % Publication open-source
            \coordinate (E2) at (3.5,0); % Incubation Apache Software Foundation
            \coordinate (E3) at (4.2,0); % Projet top-level Apache Software Foundation
            \coordinate (E4) at (7.5,0); % Adoption par Netflix
            \coordinate (E5) at (8,0); % Adoption par Spotify

            %% Filled regions
            \fill[color=facebook]
                (E1) rectangle (E2 |- 1, 1); % Facebook (2008-2009)
            \onslide<2->{\fill[color=asf]
                (E2 |- 1, 1) rectangle (F);} % Apache Software Foundation (2009-today)

            % Annotations (above)
            \draw[<-,thick,color=black] ($(E1)+(0,1.1)$) -- ($(E1)+(0,2.5)$) node [above=0pt,align=center,black] {\textbf{Septembre 2008}\\Publication\\open-source};
            \onslide<2->{\draw[<-,thick,color=black] ($(E2)+(0,1.1)$) -- ($(E2)+(0,4.5)$) node [above=0pt,align=center,black] {\textbf{Mars 2009}\\Incubation Apache\\Software Foundation (ASF)};}
            \onslide<3->{\draw[<-,thick,color=black] ($(E3)+(0,1.1)$) -- ($(E3)+(0,2.5)$) node [above=0pt,xshift=12pt,align=center,black] {\textbf{Février 2010}\\Projet ASF\\\textit{top-level}};}
            \onslide<4->{\draw[<-,thick,color=black] ($(E4)+(0,1.1)$) -- ($(E4)+(0,2)$) node [above=0pt,xshift=-12pt,align=center,black] {\textbf{2014}\\Adoption\\par Netflix};}
            \onslide<4->{\draw[<-,thick,color=black] ($(E5)+(0,1.1)$) -- ($(E5)+(0,4.2)$) node [above=0pt,align=center,black] {\textbf{2015}\\Adoption\\par Spotify};}
            % Annotation (below)
            \onslide<1>{\node[color=black,text width=13cm] at ($(E2)+(3,-2.5)$) {Cassandra, créé par Avinash Lakshman (co-auteur d’Amazon Dynamo) et Prashant Malik, sert initialement à indexer et rechercher les messages des utilisateurs de Facebook.};}

            % Annotations for the filled regions
            \draw [decorate,align=center,color=white]($(E1)+(0.85,0.1)$) -- ($(E1)+(0.8,0.1)$) node [black,midway,above=6pt] {\textcolor{white}{Facebook}};
            \onslide<2->{\draw [decorate,align=center,color=white]($(E2)+(5,0)$) -- ($(E2)+(5,0)$) node [black,midway,above=6pt] {\textcolor{white}{Apache Software Foundation}};}

            % Logos
            \node[minimum width=3cm, minimum height=2cm] at ($(E1)+(0.85,-1.1)$) {\includesvg[height=0.8cm]{img/logos/logo-facebook.svg}};
            \onslide<2->{\node[minimum width=3cm, minimum height=2cm] at ($(E2)+(5,-1.1)$) {\includesvg[height=0.8cm]{img/logos/logo-asf.svg}};}
            \onslide<4->{\node[minimum width=3cm, minimum height=2cm] at ($(E4)+(-0.4,3.8)$) {\includegraphics[height=1cm]{img/logos/logo-netflix.png}};}
            \onslide<4->{\node[minimum width=3cm, minimum height=2cm] at ($(E5)+(1.3,4.8)$) {\includegraphics[height=0.8cm]{img/logos/logo-spotify.png}};}

            %% Arrow
            \draw[->] (O) -- (F);

            %% Ticks
            \foreach \x in {0,4,...,12}
                \draw(\x cm,3pt) -- (\x cm,-3pt);

            %% Labels
            \foreach \i \j in {0/2005,4/2010,8/2015,12/2020}{
                \draw (\i,0) node[below=3pt] {\j} ;
            }
        \end{tikzpicture}
        }
    \end{figure}

    \note{
        Rappeler l'origine de Cassandra chez Facebook pour la recherche dans
        les messages, puis son passage à la fondation Apache. Les adoptions
        par Netflix et Spotify illustrent des charges distribuées à grande échelle.
        La frise permet surtout de montrer que le projet est devenu un système
        généraliste, open source et largement adopté.
    }
\end{frame}

% --------------------------------------------------------------------------------------

\begin{frame}
    Par exemple, avec Cassandra, le choix de la partition key est particulièrement important : Cassandra utilise le hachage de cette clé pour déterminer où stocker la ligne dans le cluster.

    \note{
        Faire une pause sur cette idée : Cassandra ne répartit pas les données
        au hasard du point de vue de l'application, elle applique une fonction
        de hachage à la clé de partition. La clé doit donc éviter les valeurs
        trop concentrées et correspondre aux requêtes fréquentes.
    }
\end{frame}

\begin{frame}{Un système NoSQL évolutif et distribué}

    Apache Cassandra est un système NoSQL basé sur l'\textbf{architecture BigTable de Google}, héritant ainsi du modèle de stockage orienté colonnes. Cependant, il a évolué vers un \textbf{modèle relationnel étendu}, offrant une \textbf{flexibilité bien supérieure} à ses origines.

    Pour l'interrogation des données, il propose \textbf{CQL} (Cassandra Query Language), un langage de définition de schéma et de requêtes \textbf{similaire au SQL}.

    Cassandra a été conçue dès sa création comme un \textbf{système scalable et distribué}, permettant une \textbf{scalabilité horizontale} et une \textbf{haute disponibilité}.

    \note{
        Cassandra propose CQL, qui ressemble au SQL, mais cela ne signifie pas
        qu'il s'agit d'un SGBD relationnel classique. Les tables et requêtes
        sont conçues pour la distribution et les accès connus à l'avance.
        Mettre en avant la scalabilité horizontale et l'absence de point unique
        de défaillance comme objectifs d'architecture.
    }

\end{frame}

\begin{frame}{Études de cas}
    https://cassandra.apache.org/_/case-studies.html

    Ebay : https://www.slideshare.net/slideshow/cassandra-at-ebay-13920376/13920376

    \note{
        Utiliser cette slide pour présenter plusieurs secteurs avant de détailler
        les cas Netflix et Spotify. Pour chaque cas, chercher la charge, le volume,
        la contrainte de disponibilité et la raison du choix de Cassandra.
    }
\end{frame}

\begin{frame}{Étude de cas (1)}
    https://cassandra.apache.org/_/case-studies.html

    Netflix et Cassandra (scalabilité pour le streaming).

    https://netflixtechblog.com/building-netflixs-distributed-tracing-infrastructure-bb856c319304

    Problématique.
    Solution NoSQL choisie.
    Architecture mise en place.
    Résultats (performances, coûts).

    \note{
        Netflix illustre le besoin de répartir les données et les requêtes
        sur de nombreux nœuds pour un service mondial. Insister sur le lien
        entre haute disponibilité, réplication et expérience utilisateur.
        Demander quels indicateurs permettraient de mesurer le succès :
        latence, débit, disponibilité et coût d'exploitation.
    }

\end{frame}

\begin{frame}{Étude de cas (2)}

    Netflix et Spotify (scalabilité pour le streaming).

    https://engineering.atspotify.com/2015/1/personalization-at-spotify-using-cassandra

    Problématique.
    Solution NoSQL choisie.
    Architecture mise en place.
    Résultats (performances, coûts).

    \note{
        Spotify permet d'aborder les recommandations et la personnalisation,
        avec beaucoup de lectures et un grand nombre d'utilisateurs.
        Le choix de Cassandra doit être justifié par les requêtes et le débit,
        pas seulement par le volume de données.
    }

\end{frame}
